\TemplateHeading
\templatefourrow{12.965}{课程名称}{计算机网络课程设计}{实验课时}{\experimenthours}
\templatefourrow{13.688}{实验项目}{DHCP配置与协议分析}{实验时间}{2026.11.24}
\templaterow{实验目的}{26}{
配置DHCPv4服务器和客户端，验证动态地址分配。根据报文的类型、地址和选项说明地址获取与续租的差异，辨认租期、续租时间T1和重绑定时间T2。
}
\templaterow{实验环境}{39}{
A部分使用Cisco Packet Tracer 9.0.1、Server-PT、2960交换机和PC-PT。B部分在QEMU隔离虚拟网络中运行Alpine Linux的DHCP客户端，通过Wireshark分析实捕以太网报文。虚拟接口承载释放、重新获取和续租操作。

指导教师为\teachername，课序号为06，班级为\classfullname。
}
\templaterow{\experimentlabel}{96.977}{
\subheading{实验步骤与记录方法}
1. B1启动虚拟链路抓包，确认客户端已有租约，释放后查看接口，再触发重新获取。根据事务标识找出同一次获取的报文，对照ACK和客户端上实际安装的IPv4配置。保持接口连接及有效租约。

2. A1搭建服务器、交换机及前两台PC组成的局域网。规划192.168.7.0/24，服务器使用静态地址192.168.7.2，DHCP池从192.168.7.100开始，可分配10个地址。网关和DNS保持未指定。客户端自动获取地址后分别Ping服务器。

3. A2进入Simulation模式，新增PC3并接入同一局域网。在事件列表中查看其第一次获取地址的关键事件，区分DHCP报文类型和同一报文经交换机转发的事件。记录最终IPv4配置和Ping输出。

4. B2回到保持有效租约的虚拟机，在持续抓包的条件下触发主动续租。检查续租的请求与应答，读取ACK的选项51、58和59，并与B1对照。
}
